% einzelne Kapitel werden hier eingebunden
\section{Einleitung}\label{einleitung}

Hier steht die Einleitung der Arbeit... Lorem ipsum dolor sit amet, consetetur sadipscing elitr, sed diam nonumy eirmod tempor invidunt ut labore et dolore magna aliquyam erat, sed diam voluptua. At vero eos et accusam et justo duo dolores et ea rebum. Stet clita kasd gubergren, no sea takimata sanctus est Lorem ipsum dolor sit amet. Lorem ipsum dolor sit amet, consetetur sadipscing elitr, sed diam nonumy eirmod tempor invidunt ut labore et dolore magna aliquyam erat, sed diam voluptua. At vero eos et accusam et justo duo dolores et ea rebum. Stet clita kasd gubergren, no sea takimata sanctus est Lorem ipsum dolor sit amet.

\newpage

\section{Hauptabschnitt}\label{hauptabschnitt}

Text des ersten Abschnitts... Lorem ipsum dolor sit amet, consetetur sadipscing elitr, sed diam nonumy eirmod tempor invidunt ut labore et dolore magna aliquyam erat, sed diam voluptua. At vero eos et accusam et justo duo dolores et ea rebum. Stet clita kasd gubergren, no sea takimata sanctus est Lorem ipsum dolor sit amet. Lorem ipsum dolor sit amet, consetetur sadipscing elitr, sed diam nonumy eirmod tempor invidunt ut labore et dolore magna aliquyam erat, sed diam voluptua. At vero eos et accusam et justo duo dolores et ea rebum. Stet clita kasd gubergren, no sea takimata sanctus est Lorem ipsum dolor sit amet.

\subsection{Unterabschnitt 1}\label{unterabschnitt_1}

Erstes Unterabschnitt

\subsection{Unterabschnitt 2}\label{unterabschnitt_2}

Zweites Unterabschnitt

\subsection{Unterabschnitt 3}\label{unterabschnitt_3}

Drittes Unterabschnitt

\newpage

\section{Weiterer Hauptabschnitt}\label{hauptabschnitt_2}

Text des zweiten Abschnitts... Lorem ipsum dolor sit amet, consetetur sadipscing elitr, sed diam nonumy eirmod tempor invidunt ut labore et dolore magna aliquyam erat, sed diam voluptua. At vero eos et accusam et justo duo dolores et ea rebum. Stet clita kasd gubergren, no sea takimata sanctus est Lorem ipsum dolor sit amet. Lorem ipsum dolor sit amet, consetetur sadipscing elitr, sed diam nonumy eirmod tempor invidunt ut labore et dolore magna aliquyam erat, sed diam voluptua. At vero eos et accusam et justo duo dolores et ea rebum. Stet clita kasd gubergren, no sea takimata sanctus est Lorem ipsum dolor sit amet.

\newpage

% hier können weitere Kapitel angelegt und eingetragen werden
% ....

\section{Ausblick}\label{ausblick}

Text des Ausblicks - sofern dies in der Arbeit gewünscht ist... Lorem ipsum dolor sit amet, consetetur sadipscing elitr, sed diam nonumy eirmod tempor invidunt ut labore et dolore magna aliquyam erat, sed diam voluptua. At vero eos et accusam et justo duo dolores et ea rebum. Stet clita kasd gubergren, no sea takimata sanctus est Lorem ipsum dolor sit amet. Lorem ipsum dolor sit amet, consetetur sadipscing elitr, sed diam nonumy eirmod tempor invidunt ut labore et dolore magna aliquyam erat, sed diam voluptua. At vero eos et accusam et justo duo dolores et ea rebum. Stet clita kasd gubergren, no sea takimata sanctus est Lorem ipsum dolor sit amet.

\newpage

\section{Fazit}\label{fazit}

Text des Fazits... Lorem ipsum dolor sit amet, consetetur sadipscing elitr, sed diam nonumy eirmod tempor invidunt ut labore et dolore magna aliquyam erat, sed diam voluptua. At vero eos et accusam et justo duo dolores et ea rebum. Stet clita kasd gubergren, no sea takimata sanctus est Lorem ipsum dolor sit amet. Lorem ipsum dolor sit amet, consetetur sadipscing elitr, sed diam nonumy eirmod tempor invidunt ut labore et dolore magna aliquyam erat, sed diam voluptua. At vero eos et accusam et justo duo dolores et ea rebum. Stet clita kasd gubergren, no sea takimata sanctus est Lorem ipsum dolor sit amet.


% Beispiel für Bild mit Fußnote
\begin{figure}[htb]
 \centering
 \includegraphics[width=0.4\textwidth,angle=45]{abb/logo1}
 \caption[Beispiel einer Bildbeschreibung]{Beispiel einer Bildbeschreibung\footnotemark}
\label{fig:beispiel1}
\end{figure}
\footnotetext{Bildquelle: Beispiel einer Bildquelle}

% Beispiel für Bildintegration
\begin{figure}[htb]
 \centering
 \includegraphics[width=0.3\textwidth,angle=0]{abb/logo2}
 \caption[Beschreibung]{Beschreibung}
\label{fig:Beschreibung}
\end{figure}

% Beispiel: Referenz auf Abbildung
Abbildung~\ref{fig:Beschreibung} [S.\pageref{fig:Beschreibung}]

% Beispiel: Tabelle 
\begin{center}
  \begin{tabular}{ | l | c | }
    \hline
    Überschrift 1 & Überschrift 2 \\ \hline \hline
    Info 1 & Info 2 \\ \hline
    Info 3 & Info 4 \\ \hline
    \hline
    \multicolumn{2}{|c|}{Info in einer Zelle} \\
    \hline
  \end{tabular}
\end{center}


% Beispiel für Quellcode Listings
\lstset{language=xml}
\begin{lstlisting}[frame=htrbl, caption={Die Datei {\normalfont \ttfamily  data-config.xml} dient als Beispiel für XML Quellcode}, label={lst:dataconfigxml}]
<dataConfig>
  <dataSource type="JdbcDataSource" 
              driver="com.mysql.jdbc.Driver"
              url="jdbc:mysql://localhost/bms_db"
              user="root" 
              password=""/>
  <document>
    <entity name="id"
        query="select id, htmlBody, sentDate, sentFrom, subject, textBody
        from mail">
    <field column="id" name="id"/>
    <field column="htmlBody" name="text"/>
    <field column="sentDate" name="sentDate"/>
    <field column="sentFrom" name="sentFrom"/>
    <field column="subject"  name="subject"/>
    <field column="textBody" name="text"/>
    </entity>
  </document>
</dataConfig>
\end{lstlisting}

\lstset{language=java}
\begin{lstlisting}[frame=htrbl, caption={Das Listing zeigt Java Quellcode}, label={lst:result2}]
/* generate TagCloud */
Cloud cloud = new Cloud();
cloud.setMaxWeight(_maxSizeOfText);
cloud.setMinWeight(_minSizeOfText);
cloud.setTagCase(Case.LOWER);
	    
/* evaluate context and find additional stopwords */
String query = getContextQuery(_context);
List<String> contextStoplist = new ArrayList<String>();
contextStoplist = getStopwordsFromDB(query);
	    
/* append context stoplist */
while(contextStoplist != null && !contextStoplist.isEmpty())
  _stoplist.add(contextStoplist.remove(0));
	    
/* add cloud filters */
if (_stoplist != null) {
  DictionaryFilter df = new DictionaryFilter(_stoplist);
  cloud.addInputFilter(df);
}
/* remove empty tags */
NonNullFilter<Tag> nnf = new NonNullFilter<Tag>();
cloud.addInputFilter(nnf);

/* set minimum tag length */
MinLengthFilter mlf = new MinLengthFilter(_minTagLength);
cloud.addInputFilter(mlf);

/* add taglist to tagcloud */
cloud.addText(_taglist);

/* set number of shown tags */	    
cloud.setMaxTagsToDisplay(_tagsToDisplay);
\end{lstlisting}


% Beispiel für Formeln
Die Zuordnung aller möglichen Werte, welche eine Zufallsvariable annehmen kann nennt man \emph{Verteilungsfunktion} von $X$.

\begin{quotation}
Die Funktion F: $\mathbb{R} \rightarrow$ [0,1] mit $F(t) = P (X \le t)$ heißt Verteilungsfunktion von $X$.\footnote{Mustermann, vgl.~\cite{mm2009}~[S.55]}
\end{quotation}

\begin{quotation}
Für eine stetige Zufallsvariable $X: \Omega \rightarrow \mathbb{R}$ heißt eine integrierbare, nichtnegative reelle Funktion $w: \mathbb{R} \rightarrow \mathbb{R}$ mit $F(x) = P(X \le x) = \int_{-\infty}^{x} w(t)dt$ die \emph{Dichte} oder \emph{Wahrscheinlichkeitsdichte} der Zufallsvariablen $X$.\footnote{Mustermann, vgl.~\cite{mf2005}~[S.56]}
\end{quotation}


% einfacher Zeilenabstand
\singlespacing
% Literaturliste soll im Inhaltsverzeichnis auftauchen
\newpage
\addcontentsline{toc}{section}{Literaturverzeichnis}
% Literaturverzeichnis anzeigen
\renewcommand\refname{Literaturverzeichnis}
\bibliography{Hauptdatei}

%% Index soll Stichwortverzeichnis heissen
% \newpage
% % Stichwortverzeichnis soll im Inhaltsverzeichnis auftauchen
% \addcontentsline{toc}{section}{Stichwortverzeichnis}
% \renewcommand{\indexname}{Stichwortverzeichnis}
% % Stichwortverzeichnis endgültig anzeigen
% \printindex

\onehalfspacing
% evtl. Anhang
\newpage
\addcontentsline{toc}{section}{Anhang}
\fancyhead[L]{Anhang} %Kopfzeile links
\subsection*{Anhang}\label{anhang}

Der Anhang bestehend aus Bildern und Texten...

% Beispiel für Bildintegration
\begin{figure}[htb]
 \centering
 \includegraphics[width=0.3\textwidth,angle=0]{abb/logo1}
 \caption[Abbildung im Anhang]{Abbildung im Anhang}
\label{fig:Abbildung im Anhang}
\end{figure}

Lorem ipsum dolor sit amet, consetetur sadipscing elitr, sed diam nonumy eirmod tempor invidunt ut labore et dolore magna aliquyam erat, sed diam voluptua. At vero eos et accusam et justo duo dolores et ea rebum. Stet clita kasd gubergren, no sea takimata sanctus est Lorem ipsum dolor sit amet. Lorem ipsum dolor sit amet, consetetur sadipscing elitr, sed diam nonumy eirmod tempor invidunt ut labore et dolore magna aliquyam erat, sed diam voluptua. At vero eos et accusam et justo duo dolores et ea rebum. Stet clita kasd gubergren, no sea takimata sanctus est Lorem ipsum dolor sit amet.


% Eidesstattliche Erklärung
\newpage
\addcontentsline{toc}{section}{Eidesstattliche Erklärung}
\section*{Eidesstattliche Erklärung}
\thispagestyle{empty}

\begin{verbatim}

\end{verbatim}

\begin{LARGE}Eidesstattliche Erklärung zur <-Arbeit>\end{LARGE}
\begin{verbatim}


\end{verbatim}
Ich versichere, die von mir vorgelegte Arbeit selbstständig verfasst zu haben. Alle Stellen, die wörtlich oder sinngemäß aus veröffentlichten oder nicht veröffentlichten Arbeiten anderer entnommen sind, habe ich als entnommen kenntlich gemacht. Sämtliche Quellen und Hilfsmittel, die ich für die Arbeit benutzt habe, sind angegeben. Die Arbeit hat mit gleichem Inhalt bzw. in wesentlichen Teilen noch keiner anderen Prüfungsbehörde vorgelegen.



\begin{displaymath}
% use packages: array
\begin{array}{ll}
Unterschrift:~~~~~~~~~~~~~~~~~~~~~~~~~~~~~~~~~~~~~~~~~~
& Ort, Datum:~~~~~~~~~~~~~~~~~~~~~~~~~~~~~~~~~~~~~~~~~~
\end{array}
\end{displaymath}

